\section{Introduction}
\label{sec:intro}

When a tornado or severe storm leads to a presidential disaster declaration, the federal government makes money available not only for recovery but also for reducing the damage from the next event. The main channel is the Federal Emergency Management Agency's (FEMA) Hazard Mitigation Grant Program (HMGP). Its funding is triggered by the declaration and sized as a share of the federal assistance that the disaster generates \parencite{horn2022crs}. Between 1996 and 2019 the program awarded \$22.4 billion, far more than FEMA's pre-disaster mitigation programs \parencite{horn2022crs}. Studies of FEMA mitigation spending report benefits well above costs \parencite{rose2007bca,godschalk2009foresight} and measurable reductions in later losses \parencite{davlasheridze2017adaptation,miao2025loss}. Where this money is spent therefore affects which places carry less residual risk, and lower expected losses, into the next storm.

The money does not flow automatically to the places a disaster struck. HMGP is an intergovernmental grant with a state intermediary: funds go to the state, which selects among subapplications prepared by counties, cities, school districts, special districts, and other eligible entities \parencite{horn2022crs}. Each subapplicant must hold a FEMA-approved local mitigation plan, show that the project is cost-effective, and usually provide a 25 percent nonfederal match \parencite{horn2022crs,frimpong2023federal}. A declaration thus creates an opportunity for federally financed public investment, but the opportunity becomes a project only if a local government decides to apply, scopes a project, and prepares the documentation. Counties that share the same disaster declaration can therefore end up with very different levels of mitigation investment. Political incentives already favor spending on relief over spending on preparedness \parencite{healy2009myopic}, and the distributional consequences of FEMA's risk-reduction programs are a recurring concern \parencite{noonan2018flood,miao2025equity}.

Small jurisdictions have reasons to apply less often. Much of the work of preparing an application does not shrink with the size of the place: an engineering estimate, a benefit--cost analysis, environmental review, and the administration of a grant require similar specialized effort whether the project serves a county of a few thousand residents or one of several hundred thousand. When these costs are largely fixed, a small government pays more per resident to obtain the same type of project, and it is less likely to employ staff who prepare such applications regularly. Research on intergovernmental grants finds that the localities with the greatest needs are often the least equipped to compete for funds \parencite{collins2006redistributive,hall2008grant}. In hazard mitigation, rural local governments with lower measured capacity secure less competitive federal resilience funding \parencite{smith2025rural}, FEMA programs that are administered through state and local governments tend to favor better-resourced counties \parencite{miao2025equity}, property buyouts have concentrated in wealthier and more populous counties \parencite{mach2019buyouts}, and rural counties have received less HMGP assistance per recipient than urban counties \parencite{seong2022naturally}. This fixed-cost reasoning motivates the analysis. Our design does not measure application costs directly, and we do not treat the reasoning as a tested result.

Three features of the existing evidence limit what it can say about small places. First, most studies observe only awards, so a shortfall in funding cannot be attributed to fewer applications, more rejections, or lower need. Second, comparisons across counties hit by different disasters, in different states and years, mix differences in local governments with differences in hazard severity, state administration, and the size of the HMGP allocation. Third, wind hazards have received less attention than floods, although tornado safe rooms and wind retrofits are a large share of HMGP construction and have long been studied for their cost-effectiveness \parencite{simmons2006shelters}.

This article asks whether, and at which stage, small counties fall behind in turning tornado and severe-storm disaster declarations into funded mitigation projects. We assemble every county designated in a tornado or severe-storm major disaster declaration from 1999 to 2021 and link each county--declaration to the HMGP subapplications and projects filed under that declaration, including requests that were later withdrawn or denied. Tornado exposure is measured from Storm Prediction Center (SPC) tracks that crossed populated block groups, so that exposure does not depend on damage reports. Counts and finances of local governments come from the Census of Governments. The central comparison is between counties designated in the same declaration, which share the incident, the state administration, and the HMGP allocation.

Within the same declaration, the smallest third of counties are 8.8 percentage points less likely to obtain any HMGP project, against a mean of 27.1 percent among larger counties, and about half as likely to obtain a safe room or wind retrofit. Participation falls steadily with population across the whole distribution rather than only among the very smallest counties. The shortfall arises before any request is reviewed: conditional on applying, small counties have their requests approved, and withdraw them, at rates similar to those of larger counties. About two-fifths of the gap is associated with the smaller number of municipalities, school districts, and special districts that could apply in small counties. A gap remains for the county government itself, of which every county has exactly one. When the same county is designated in several declarations, a tornado that crosses a populated place is followed by more applications in larger counties, whereas the estimated response of small counties is negative and imprecise; the difference between the two is marginally significant. Per-resident fiscal measures do not account for the gap, a 2018 federal provision that reimburses subrecipients' management costs did not detectably narrow it, and the gap is smaller in states whose agencies historically filed a larger share of applications.

The article contributes to the regional economics of intergovernmental transfers in three ways. First, it locates a geographic disparity in access to federally funded public investment at a specific administrative margin, the decision to apply, using application records that include withdrawn and denied requests. Second, it separates a component of the county-size gap that reflects the institutional structure of local government, namely the number of potential applicants, from a residual that remains for a single, comparable applicant in every county. Third, it examines several explanations proposed in the literature, including fiscal constraints, state delivery of applications, and federal reimbursement of administrative costs, within a common comparison. The evidence is observational: it shows where and for which places the gap arises, not the causal effect of local administrative capacity. For policy, the results suggest that disparities in grant-funded investment across places can arise even when funding formulas respond to hazards, because access depends on whether a local government applies.

Section~\ref{sec:background} describes the program and related research. Section~\ref{sec:data} presents the data and empirical approach, Section~\ref{sec:results} the results, and Section~\ref{sec:discussion} their interpretation, policy implications, and limitations. Section~\ref{sec:conclusion} concludes.
